% ============================================================
% Capítulo 4: IA como Parceira de Engenharia
% ============================================================
\chapter{IA como Parceira de Engenharia}

\section{Introdução}

% [Placeholder: O novo paradigma]
% - A IA não vai substituir engenheiros, mas engenheiros que usam IA
%   vão substituir os que não usam
% - Mudança de mentalidade: de operador a supervisor
% - O que a IA faz bem e o que ela não faz

\section{O Que São LLMs e Como Funcionam}

% [Placeholder: Fundamentos técnicos acessíveis]
% - Treinamento e predição de tokens
% - Limitações intrínsecas (alucinação, contexto, conhecimento)
% - Por que LLMs não "pensam" — elas simulam padrões
% - Implicações para geração de código

\section{Modelos de Colaboração Humano-IA}

% [Placeholder: Formas de trabalhar com IA]
% - IA como par programador
% - IA como revisor de código
% - IA como gerador de boilerplate
% - IA como tutor e explicador
% - IA como explorador de alternativas

\section{Limitações e Armadilhas}

% [Placeholder: O que a IA não faz bem]
% - Código que compila mas não resolve o problema
% - Falta de compreensão do domínio
% - Vieses do training data
% - Segurança e vazamento de informações sensíveis
% - Dependência excessiva e perda de pensamento crítico

\section{O Papel do Julgamento Humano}

% [Placeholder: Onde o engenheiro é insubstituível]
% - Compreensão do contexto de negócio
% - Decisões de trade-off
% - Avaliação de impacto sistêmico
% - Responsabilidade final pelo código

\section{Ferramentas e Ecossistema}

% [Placeholder: Panorama de ferramentas]
% - GitHub Copilot e similares
% - Agents autônomos
% - IDEs com integração de IA
% - Como escolher ferramentas

\section{Workflow com IA}

% [Placeholder: Integrando IA no dia a dia]
% - Quando usar IA e quando não usar
% - Estrutura de uma sessão de trabalho com IA
% - Validação sistemática do código gerado
% - Mantendo o controle do design
% - Conexão com Capítulo 1: TDD como guardrail da IA
% - Spec Driven Development + RM-ODP como contexto para prompts (Cap. 1)

\section{Conclusão}

% [Placeholder: Resumo]
% - IA é uma ferramenta poderosa, não uma solução mágica
% - O pensamento de engenharia é mais importante que nunca

\section*{Exercícios}

% [Placeholder: Exercícios]
% 1. Use uma ferramenta de IA para gerar código para um problema
%    simples. Analise criticamente o resultado.
% 2. Liste 5 situações em que você NÃO usaria IA para gerar código.
% 3. Crie um checklist de validação para código gerado por IA.
